\documentclass[11pt]{article}

\usepackage[margin=25mm]{geometry}

\usepackage{amsmath,amssymb,fontspec,unicode-math}

\IfFontExistsTF{DejaVu Serif}{\setmainfont{DejaVu Serif}}{\setmainfont{Latin Modern Roman}}

\IfFontExistsTF{Latin Modern Math}{\setmathfont{Latin Modern Math}}{\setmathfont{DejaVu Math TeX Gyre}}

\usepackage{xurl,hyperref}

\setlength{\parindent}{0pt}\setlength{\parskip}{6pt}

\begin{document}

\title{Decomposing Interventional Causality into Synergistic, Redundant, and Unique Components}

\author{ELVIA preliminary AI analysis}\date{21 October 2026}\maketitle

Preliminary assessment: Major.\par

Independent scientific review remains pending. Author adoption has not been requested.

A fully supported binary causal model makes MACE decrease when interventions enlarge, refuting unrestricted monotonicity/nonnegativity claims. Plotting can then omit or clip signed contributions.

\section*{Review scope}

Rapid current-version reassessment of every substantive concern in the supplied LaTeX report: the implicated paper passages and relevant author-released source paths were inspected, existing proposed corrections were briefly checked, and impact was calibrated to the demonstrated scope. No author-provided or historical verification program was executed.

\section*{Problem 1: Enlarging the intervention set can reduce MACE, invalidating the unrestricted nonnegativity theorem}

Source: Section 3.2, Lemma 1 and Theorem 2, PDF p. 5; proof, PDF p. 6; Appendix A.1, Lemma 2, PDF p. 13.

Assessment: supported; impact: major

\subsection*{Source claim (paraphrase)}

Paraphrase of the complete relevant claims: Appendix A.1, Lemma 2 (p. 13), asserts that \(A\subseteq B\) implies \(m(A)\leq m(B)\). Section 3.2, Lemma 1 (p. 5), asserts monotonicity of \(r\) on the redundancy order. Theorem 2 (p. 5; proof on p. 6) then asserts \(C(\alpha;Y)\geq0\) for all antichains. No assumption of independent source variables or monotonicity of the actual interventional response ranges is attached to these claims. The paper explicitly discusses correlated dynamical contexts in Section 4.2.

\subsection*{Analysis}

Let \(A,H\) be independent fair bits, \(E\sim\operatorname{Bernoulli}(1/4)\) independent, \(B=H\oplus E\), \(C=H\), and \(Y=A\oplus B\oplus C\). All observable source assignments have positive probability and no source causes another. Under \(\operatorname{do}(A)\), \(Y=A\oplus E\), giving \(\operatorname{MACE}(A)=1/2\). Under \(\operatorname{do}(A,B)\), \(C\) remains fair and \(Y\) is fair for every imposed assignment, giving \(\operatorname{MACE}(AB)=0\).

Adding interventions replaces mechanisms; the resulting response distributions need not contain those of the smaller intervention family. The decrease contradicts unrestricted intervention monotonicity, and the stated inversion produces a negative partial contribution \(-1/2\). Correlated contexts are permitted by the source, so independence cannot be silently added to remove the example.

For an explicit two-source decomposition, choose analyzed sources \(\{A,B\}\) while \(C\) remains background. Then \(m(A)=1/2\), \(m(B)=0\), \(m(AB)=0\) and the minimum redundancy is \(0\). Möbius reconstruction gives the joint-only atom \(m(AB)-\max\{m(A),m(B)\}=-1/2\). Every source assignment is supported, so this does not rely on unspecified off-support responses.

\subsection*{Finding and scope}

The unrestricted monotonicity and universal nonnegativity claims are false. The decomposition can still be studied with an explicit monotonicity assumption or a separately defined envelope; those changes alter its scope or meaning.

\subsection*{Proposed correction}

State and justify a monotonicity assumption for the intervention-effect set function, or define a monotone-envelope variant and disclose its changed scientific meaning. Neither is a replacement proof of the unrestricted original claim.

A local correction or qualification is proposed where specified; author adoption is not established.

\section*{Problem 2: The cellular-automaton plot omits signed contributions produced by the decomposition}

Source: Section 4.2, PDF p. 7; Figure 3, PDF p. 8; released notebook cells 5–6.

Assessment: supported; impact: minor

\subsection*{Source claim (paraphrase)}

Section 4.2, PDF p. 7, and Figure 3, PDF p. 8, compare cellular-automaton contexts. The authors' notebook, zero-based cells \(5\)–\(6\), implements simulation and plotting with a positive-only atom-selection threshold and nonnegative display limits.

\subsection*{Analysis}

Atoms never exceeding \(0.01\) are omitted by the positive-only mask, while display limits \([0,1.1]\) clip negative contributions of retained atoms. Once signed contributions are possible, those operations can hide them. The released settings of \(1000\) steps with \(0.05\) burn-in also differ from the stated \(10000\) steps with the first \(500\) removed. No new cellular-automaton simulation or negative-context frequency is claimed.

\subsection*{Finding and scope}

This is a supported local finding within the stated source/formulation scope. Its effect on the paper's aggregate experimental results remains unmeasured.

\section*{Limits}

No author-provided program, training procedure, model inference or wet-lab experiment was executed in this review.\par

Historical execution results are retained only with that attribution; a source-level defect alone does not identify the program or configuration used for the published experiments.\par

No author contact, adoption, independent human scientific review or novelty determination is established.\par

The findings apply only to the identified claims, source versions and stated assumptions; they do not certify or invalidate the entire paper.\par

\section*{Source identity}

NeurIPS 2025 official proceedings PDF, freshly retrieved from the current publisher landing page

\url{https://papers.nips.cc/paper_files/paper/2025/file/535baca36883a4e0450423b26b150a48-Paper-Conference.pdf}

Reviewed PDF SHA-256: \nolinkurl{0d227dd1ac0a5eff8df75e284daded1ea911514c3a73fe24f1f790d7a892c3b9}

The exact inspected current paper PDF is identified by its SHA-256. The current publisher PDF was freshly downloaded and is byte-identical to the previously retained proceedings copy. Latest author-manuscript history checked: 2501.11447v2 (last updated 2025-09-19T10:54:57Z). The current formal proceedings remain the primary manuscript; the latest indexed preprint does not postdate or replace the formal work in this channel. Current author resource: https://github.com/AJnsm/causalDecomposition at commit d3d6fad3349461f000a1052861fb74fd47111008.  An exact PDF acquisition timestamp is not established in this reassessment; source-currentness checks have their own recorded timestamps and file modification times were not substituted.

\end{document}
